%%%%%%%%%%%%%%%%%%%%%%%%%%%%%%%%%%%%%%%%%
% FRI Data Science_report LaTeX Template
% Version 1.0 (28/1/2020)
% 
% Jure Demšar (jure.demsar@fri.uni-lj.si)
%
% Based on MicromouseSymp article template by:
% Mathias Legrand (legrand.mathias@gmail.com) 
% With extensive modifications by:
% Antonio Valente (antonio.luis.valente@gmail.com)
%
% License:
% CC BY-NC-SA 3.0 (http://creativecommons.org/licenses/by-nc-sa/3.0/)
%
%%%%%%%%%%%%%%%%%%%%%%%%%%%%%%%%%%%%%%%%%


%----------------------------------------------------------------------------------------
%	PACKAGES AND OTHER DOCUMENT CONFIGURATIONS
%----------------------------------------------------------------------------------------
\documentclass[fleqn,moreauthors,10pt]{ds_report}
\usepackage[english]{babel}

\graphicspath{{fig/}}




%----------------------------------------------------------------------------------------
%	ARTICLE INFORMATION
%----------------------------------------------------------------------------------------

% Header
\JournalInfo{FRI Natural language processing course 2026}

% Interim or final report
\Archive{Project report} 
%\Archive{Final report} 

% Article title
\PaperTitle{Redacted} 

% Authors (student competitors) and their info
\Authors{Miha Mitič in Jošt Eržen}

% Advisors
\affiliation{\textit{Advisors: Slavko Žitnik}}

% Keywords
\Keywords{PC, hardware, computers, domain specific AI, assistant}
\newcommand{\keywordname}{Keywords}


%----------------------------------------------------------------------------------------
%	ABSTRACT
%----------------------------------------------------------------------------------------

\Abstract{

We propose a Redacted — a domain-specific conversational AI that recommends PC components based on user needs and budget. General-purpose LLMs often provide inaccurate component advice due to stale or hallucinated specifications - we will build an assistant with explicit compatibility rules to ground recommendations in verified manufacturer data.

}

%----------------------------------------------------------------------------------------

\begin{document}

% Makes all text pages the same height
\flushbottom 

% Print the title and abstract box
\maketitle 

% Removes page numbering from the first page
\thispagestyle{empty} 

%----------------------------------------------------------------------------------------
%	ARTICLE CONTENTS
%----------------------------------------------------------------------------------------

\section*{Introduction}

Choosing PC components is a common technical challenge for newcomers to the hobby, yet general-purpose LLMs often provide incorrect information (e.g., CPU socket compatibility, number of PCIe slots, bandwidth...). Our goal is to design an AI assistant that provides trustworthy, source-grounded component recommendations for PC builders. We aim to speed up the component selection process while recommending components that would suit the consumer's needs best. 

%------------------------------------------------

\section*{Proposed Dataset}

We will compile a corpus of official product specification pages and manuals from major component manufacturers (ASUS, MSI, Gigabyte, ASRock, AMD, Intel, NVIDIA). For each component we will extract both structured attributes and unstructured text segments (datasheets, compatibility notes, BIOS guidance) to support RAG. The initial scope targets consumer-grade parts released in roughly the last three years so the assistant remains relevant for typical PC builders.



\section*{Initial approach}

We will scrape and normalize official specification pages into a structured schema plus text chunks, then enforce explicit compatibility rules (socket, chipset, RAM type, PCIe lanes, PSU headroom) to ensure each candidate build is compatible before recommendations are generated. A Retrieval-Augmented Generation layer retrieves the most relevant spec passages for each query, and an instruction-tuned LLM combines that evidence with rule verdicts to produce conversational recommendations with source citations. 

\section*{Related work}
RAG-based systems demonstrate that injecting retrieved passages before decoding substantially improves factual accuracy on knowledge-intensive tasks, yet still struggles when key specifications are missing \cite{lewis2020rag,izacard2021leveraging}. Instruction-tuned models offer better conversational control but inherit the same grounding limitations unless connected to authoritative data sources \cite{ouyang2022training}. Bengio argues for hybrid neural–symbolic approaches to ensure dependable reasoning in high-stakes settings, motivating our combination of RAG with explicit compatibility rules \cite{bengio2020reconsidering}. Finally, tools like the PCPartPicker compatibility engine show that deterministic rule bases prevent hardware mismatches but lack conversational flexibility; our assistant aims to merge their reliability with LLM-driven dialogue \cite{pcpartpicker2026compatibility}.



%----------------------------------------------------------------------------------------
%	REFERENCE LIST
%----------------------------------------------------------------------------------------
\bibliographystyle{unsrt}
\bibliography{report}


\end{document}